%% kapa: exam=1 round=29 year=2018 period=1 form=A subject=economics subject_ko=경제학원론 seq=18 number=58
%% source: Q-Net 감정평가사 기출문제 (cst00309) artlSeq=5202808 file=extracted/1008/5202808/2018년 제29회 감정평가사 1차시험 기출문제/2018년 제29회 감정평가사 1차시험 기출문제/1교시 A형.pdf
%% conversion: kordoc 4.24.51 -> tools/gen1.py (auto); review=proofread-sonnet; flags=none
\begin{question}{58}
두 재화 맥주( $B$)와 커피( $C$)를 소비하는 두 명의 소비자 1과 2가 존재하는 순수교환경제를 가정한다. 소비자 1의 효용함수는 $U_{1}(B_{1},C_{1})=\min\{B_{1},C_{1}\}$, 소비자 2의 효용함수는 $U_{2}(B_{2},C_{2})=B_{2}+C_{2}$이다. 소비자 1의 초기 부존자원은 (10,20), 소비자 2의 초기 부존자원은 (20,10)이고, 커피의 가격은 1이다. 일반균형(general equilibrium)에서 맥주의 가격은? (단, 초기 부존자원에서 앞의 숫자는 맥주의 보유량, 뒤의 숫자는 커피의 보유량이다.)
\begin{choices}
  \item $1/3$
  \item $1/2$
  \item 1
  \item 2
  \item 3
\end{choices}
\end{question}
\answer{3}
